\documentclass[11pt,a4paper]{article}

% Paper-style figure catalogue for LogoLens MSc Dissertation v17.
% The figure blocks can also be copied individually into the main manuscript.
\usepackage[T1]{fontenc}
\usepackage[utf8]{inputenc}
\usepackage{lmodern}
\usepackage{microtype}
\usepackage[a4paper,margin=25mm]{geometry}
\usepackage{graphicx}
\usepackage{caption}
\usepackage{float}
\usepackage{placeins}
\usepackage[hidelinks]{hyperref}

\graphicspath{{LogoLens_MSc_Dissertation_v17_media/}}

\captionsetup{
  font=small,
  labelfont=bf,
  labelsep=period,
  justification=justified,
  singlelinecheck=false
}

\setlength{\textfloatsep}{12pt plus 2pt minus 2pt}
\setlength{\floatsep}{12pt plus 2pt minus 2pt}
\setlength{\intextsep}{12pt plus 2pt minus 2pt}

\title{LogoLens: Paper-Style Figure Catalogue}
\author{MSc Dissertation, Version 17}
\date{}

\begin{document}
\maketitle
\listoffigures
\clearpage

% Figure 1
\begin{figure}[!htbp]
  \centering
  \includegraphics[width=0.98\linewidth]{figure_01_yolo26_rfdetr_architecture.png}
  \caption[YOLO26 and RF-DETR conceptual architecture comparison]{%
    Conceptual comparison of the evaluated YOLO26 and RF-DETR detection
    routes. The diagram shows the task-relevant prediction pathways and their
    common output schema. It is not a layer-complete network specification and
    does not imply equivalent accuracy or computational efficiency.}
  \label{fig:yolo-rfdetr-architecture}
\end{figure}
\FloatBarrier

% Figure 2
\begin{figure}[!htbp]
  \centering
  \includegraphics[width=0.96\linewidth]{figure_02_research_framework.png}
  \caption[Research framework]{%
    Research framework linking sponsorship accountability to observable
    exposure, leakage-aware computer vision and auditable visibility evidence.
    The diagram is conceptual and contains no measured experimental results.}
  \label{fig:research-framework}
\end{figure}
\FloatBarrier

% Figure 3
\begin{figure}[!htbp]
  \centering
  \includegraphics[width=0.98\linewidth]{figure_03_research_design_case_boundary.png}
  \caption[Research design and case boundary]{%
    Leakage-aware research design and case boundary. Training, validation and
    held-out matches have non-overlapping roles. The evaluated detection study
    is bounded to one club, the white home kit, broadcast and highlight footage,
    and three unseen test matches. Visibility aggregation consumes evaluated
    detections but remains beyond the validated detection boundary.}
  \label{fig:research-design-boundary}
\end{figure}
\FloatBarrier

% Figure 4
\begin{figure}[!htbp]
  \centering
  \includegraphics[width=0.94\linewidth]{figure_04_split_and_leakage.png}
  \caption[Match-disjoint split and temporal leakage]{%
    Match-disjoint dataset design and measured inflation on familiar match
    material. Complete matches, rather than neighbouring video frames, define
    the data partitions. The familiar-material score is a temporal-leakage
    diagnostic and is not treated as a second held-out test result.}
  \label{fig:split-and-leakage}
\end{figure}
\FloatBarrier

% Figure 5
\begin{figure}[!htbp]
  \centering
  \includegraphics[width=0.98\linewidth]{figure_05_system_architecture.png}
  \caption[LogoLens system architecture]{%
    LogoLens processing pipeline. The detector inset presents RF-DETR Small and
    YOLO26 as alternative backends behind a common adapter. Both return class,
    confidence, bounding-box, frame-index and timestamp records, allowing final
    reports to remain traceable to the source footage.}
  \label{fig:system-architecture}
\end{figure}
\FloatBarrier

% Figure 6
\begin{figure}[!htbp]
  \centering
  \includegraphics[width=0.96\linewidth]{figure_06_visibility_aggregation.png}
  \caption[Visibility-quality and duration aggregation]{%
    Transformation from detection records to prototype visibility evidence.
    The inset distinguishes measured geometry, model-derived confidence and
    researcher-defined transformations. Temporal segments, raw visible seconds
    and quality-weighted exposure remain subject to manual timing and human
    readability validation.}
  \label{fig:visibility-aggregation}
\end{figure}
\FloatBarrier

% Figure 7
\begin{figure}[!htbp]
  \centering
  \includegraphics[width=\linewidth]{figure_07_training_and_heldout_evidence.png}
  \caption[RF-DETR training trajectory and held-out evidence]{%
    RF-DETR Small training and held-out evidence. (a) Composite training and
    validation loss across the two retained runs joined at the resume boundary.
    (b) Regular and exponential-moving-average validation
    mAP@[0.50:0.95], with epoch 36 selected by the maximum EMA value.
    (c) Operating-point confusion matrix at confidence 0.35 and IoU 0.50;
    colour is normalised by ground-truth row and labels report raw counts.
    ``No object'' denotes misses and unmatched predictions.
    (d) Per-brand AP@0.50 with held-out support in parentheses; orange denotes
    classes with fewer than five boxes. CCH is omitted from panels (c) and (d)
    because the held-out set contains no CCH ground truth and its AP is not
    estimable.}
  \label{fig:training-heldout-evidence}
\end{figure}
\FloatBarrier

% Figure 8
\begin{figure}[!htbp]
  \centering
  \includegraphics[
    width=\linewidth,
    height=0.78\textheight,
    keepaspectratio
  ]{figure_08_occlusion_examples_v2.png}
  \caption[Qualitative held-out occlusion examples]{%
    Qualitative held-out examples using a three-level occlusion scale.
    Solid boxes denote matched ground truth, the dashed box denotes a missed
    ground-truth instance, and box colour identifies the sponsor class.
    The cases illustrate appearance conditions and detector outcomes; they are
    not an aggregate clear-versus-occluded performance benchmark.}
  \label{fig:occlusion-examples}
\end{figure}
\FloatBarrier

% Figure 9
\begin{figure}[!htbp]
  \centering
  \includegraphics[
    width=\linewidth,
    height=0.78\textheight,
    keepaspectratio
  ]{figure_09_detection_evidence.png}
  \caption[Representative held-out successes and failure modes]{%
    Representative held-out outcomes from the selected RF-DETR Small EMA
    checkpoint at epoch 36, regenerated at confidence 0.35 and IoU 0.50.
    Box colour identifies the sponsor class; solid boxes are correct
    predictions, short-dashed boxes labelled FP are false positives, and
    long-dashed boxes labelled Missed are missed ground truth. Source frames
    retain their natural colour, and the cases follow the predefined selection
    and annotation-audit rules.}
  \label{fig:detection-evidence}
\end{figure}
\FloatBarrier

% Figure 10
\begin{figure}[!htbp]
  \centering
  \includegraphics[width=\linewidth]{figure_10_visibility_timeline.png}
  \caption[Prototype brand-level visibility timeline]{%
    Prototype brand-level visibility timeline in the LogoLens reporting
    interface. Colour distinguishes sponsor classes, horizontal position
    encodes elapsed video time, and each bar represents a continuous detected
    segment. The interface demonstrates temporal reporting and traceability;
    the displayed intervals are not treated as manually validated duration
    ground truth.}
  \label{fig:visibility-timeline}
\end{figure}
\FloatBarrier

% Cross-reference examples for the main manuscript:
% Figure~\ref{fig:yolo-rfdetr-architecture}
% Figures~\ref{fig:training-heldout-evidence} and~\ref{fig:detection-evidence}

\end{document}
